\scanpage{34}
\mspage{13}
\margin{\uline{Factoriza-tion} \uline{Problem}}
\margin{$\unsure{\Phi_0}$ given, $\unsure{\Phi}$ variable. \deleted{NB we need only $\unsure{\Phi}$}}
$M$ a category, $\Phi_0\subset\mathrm{Fl}(M)$, $\Psi=(\Phi_0^{\deleted{*}})_*$, $\Phi_0\subset\Phi\subset\overline\Phi_0=\Psi^*$, we want to deduce conditions on $\Phi_0,\Phi$ insuring that the pair $(\Phi,\Psi)$ satisfies factorization condition, i.e.
\[
\forall f\in\mathrm{Fl}(M),\quad\exists\text{ factorization }f=pi,\quad i\in\Phi,\quad p\in\Psi.
\]
\margin{NB. we do not use $\Phi\subset\unsure{\overline\Phi_0|\Phi_0}$.}
The idea is to take \inserted{(for given $f$)} all diagrams
\diagram{\begin{tikzcd}
A\arrow[r,"v"]\arrow[d,"g"']&X\arrow[d,"f"]\\
B\arrow[r,"u"']&Y
\end{tikzcd}} $(D)$ with $g\in\Phi_0$
giving rise to
\diagram{\begin{tikzcd}
A\arrow[r,"v"]\arrow[d,"\unsure{a|g}"']\arrow[dr,phantom,"\mathrm{coc}"]&X\arrow[d,"\unsure{i_D|g_D}"]\arrow[dd,bend left=55,"f"]\\
B\arrow[r]&X_D=X\amalg_A B\arrow[d,"p_D"]\\
&Y
\end{tikzcd}}
and to define $\deleted{\unsure{\Sigma_1}(f)}\inserted{X_1(f)}=\deleted{\varinjlim}\coprod_{\text{all }D}(X)\,X_D$ (amalgamated sum under $X$)
so that we get factorization
\diagram{\begin{tikzcd}
X\arrow[r,"i_1(f)"]\arrow[d,"f"']&\deleted{\Sigma}X_1(f)\arrow[dl,"\Sigma_1(f)"]\\
Y&
\end{tikzcd}}
\deleted{with} (hopefully) $i_1(f)\in\Phi$.

For this first step, we need

$a)$ $\Phi_0$ is small (we'll see later how to relax this condition)

$b)$ Amalgamated sums indexed by $I$ with $\operatorname{card}I\leq\operatorname{card}\Phi_0$ exist in $M$, \deleted{and}
\inserted{rather, $\overbrace{\operatorname{card}\Phi_0\times c}^{\sigma}$, where $c$ is $\sup(\operatorname{card}\operatorname{Hom}(A,X)\times\operatorname{card}\operatorname{Hom}(B,Y))$ for $\unsure{u|g}:A\to B$ in $\Phi_0$}

$c)$ $\Phi_{\deleted{\unsure{0}}}$ stable by \inserted{\uline{Composition, and by}} cobase extensions, and \inserted{$(d_0)$} by filtering direct limits \inserted{all} indexed by a set $I$ of cardinality $\leq\operatorname{card}\Phi_0$ \inserted{\unsure{v. thus} $\leq\sigma$}.
\margin{do with \unsure{strengthened}}

We then are making a transfinite induction to define $X_\alpha(f)$ and
\diagram{\begin{tikzcd}
X\arrow[r,"i_\alpha(f)"]\arrow[d,"f"']&X_\alpha(f)\arrow[dl,"\Sigma_\alpha(f)"]\\
Y&
\end{tikzcd}}
$i_\alpha(f)\in\Phi$
with $\alpha$ any ordinal \deleted{$\unsure{\Phi}$} $<\alpha_0$ ($\alpha_0$ \deleted{first} ordinal \deleted{\unsure{of the cardinal}\illegible{2}} \inserted{number to be fixed} \deleted{\unsure{define later in terms of $\pi$}} \inserted{later}).
\[
\left\{\begin{aligned}
\Sigma_{\alpha+1}(f)&=\Sigma_1(\Sigma_\alpha(f)),\quad X_{\alpha+1}(f)=\text{source}\,\Sigma_{\alpha+1}(f),\quad i_{\alpha+1}(f)=\text{composition}\\[4em]
\Sigma_\alpha(f)&=\varinjlim_{\alpha'<\alpha}\Sigma_{\alpha'}(f)\quad\text{if $\alpha$ is a limiting ordinal.}
\end{aligned}\right.
\]
\diagram{\begin{tikzcd}
X\arrow[r,"i_\alpha(f)"]\arrow[d,"f"']&X_\alpha(f)\arrow[r,"i_1(\Sigma_\alpha(f))"]\arrow[dl,"\Sigma_\alpha(f)"]&X_{\alpha+1}(f)\arrow[dll,"\Sigma_{\alpha+1}(f)=\Sigma_1(\Sigma_\alpha(f))"]\\
Y&&
\end{tikzcd}}
\deleted{We} We assume that $\Phi$ is stable too by \inserted{the} filtering direct limits \deleted{of cardinality $\leq c$} which occur in the several transfinite induction steps.
Let $\pi$ be a cardinal such that \deleted{all objects of $\Phi_0$} \inserted{for any $g\in\Phi_0$, source and target of $g$} are $\pi$-\unsure{access. ob.}
--- i.e. $\operatorname{Hom}(A,-):M\to\mathrm{Sets}$ commutes to filtering direct limits which are \deleted{\unsure{great wr}} ``large with resp. to $\pi$''. We may take $\pi=\sup_{g\in\Phi_0}\pi_g$, where
\mspage{33}

